\documentclass[11pt]{article}

\usepackage[T1]{fontenc}
\usepackage[utf8]{inputenc}
\usepackage{amsmath,amssymb,amsthm}
\usepackage{mathtools}
\usepackage[margin=2.5cm]{geometry}
\usepackage{booktabs}
\usepackage{tabularx}
\usepackage{enumitem}
\usepackage{url}
\usepackage[colorlinks=true,linkcolor=blue,urlcolor=blue]{hyperref}

% Lean declaration names, typeset in typewriter with line breaks at dots.
\DeclareUrlCommand\lean{\urlstyle{tt}}

\newtheorem{definition}{Definition}[section]
\newtheorem{proposition}[definition]{Proposition}
\theoremstyle{remark}
\newtheorem{remark}[definition]{Remark}
\newtheorem{choice}[definition]{Choice}

\newcommand{\proet}{\mathrm{pro\acute{e}t}}
\newcommand{\Sh}{\mathrm{Sh}}
\newcommand{\Mod}{\mathrm{Mod}}
\newcommand{\Ab}{\mathrm{Ab}}
\newcommand{\Hom}{\mathrm{Hom}}
\newcommand{\RHom}{R\mathcal{H}om}
\newcommand{\D}{\mathrm{D}}
\newcommand{\Res}{\mathrm{Res}}
\newcommand{\Ext}{\mathrm{Ext}}
\newcommand{\Coext}{\mathrm{Coext}}
\newcommand{\Qlbar}{\overline{\mathbb{Q}}_\ell}
\newcommand{\cO}{\mathcal{O}}
\newcommand{\lotimes}{\otimes^{\mathrm{L}}}

\title{Pro-\'etale derived categories and six functors:\\
the constructions chosen in the formalization}
\author{Weil conjectures project}
\date{October 2026}

\begin{document}
\maketitle

\begin{abstract}
This note records the current state of the pro-\'etale part of the six-functor formalism
in the Lean formalization of the Weil conjectures, with emphasis on the
\emph{mathematical choices} made: which categories serve as derived categories, how the
$\ell$-adic coefficient rings are modelled, how constructibility is defined, and how each
of the six functors is constructed. It summarizes the work on the branches
\texttt{ring-objects-for-coefficients}, \texttt{extend-scalars-module-object},
\texttt{bounded-constructible-module-object-of-sheaf} and
\texttt{module-object-six-functors}, on top of what is already on \texttt{master}.
\end{abstract}

\tableofcontents

%======================================================================
\section{Overview}
%======================================================================

The formalization follows Bhatt--Scholze, \emph{The pro-\'etale topology for schemes},
\S6.8, for $\ell$-adic coefficients. Two parallel layers coexist.

\begin{enumerate}[label=(\Alph*)]
\item \textbf{Fixed coefficient category.} Sheaves on $X_\proet$ with values in a fixed
  closed symmetric monoidal Grothendieck abelian category $\mathcal{A}$ (the intended case
  being $\mathcal{A} = \Mod_\Lambda$ for a ring $\Lambda$ such as $\mathbb{Z}/\ell^n$).
  The six functors are built on $\Sh(X_\proet,\mathcal{A})$ and its derived category.
  This is on \texttt{master} (folder \lean{SixFunctors/ProEtale}) and is moved to
  \lean{SixFunctors/ProEtale/FiniteCoefficient} on \texttt{module-object-six-functors}.
\item \textbf{Module objects over a coefficient ring sheaf.} Sheaves of modules over a
  sheaf of rings $\mathcal{R}_X$ on $X_\proet$ which genuinely varies with $X$, for
  instance $\mathbb{Z}_{\ell,X}$, $\mathbb{Q}_{\ell,X}$, $E_X$, $\cO_{E,X}$. These are
  modelled as \emph{module objects} over a \emph{ring object} of the category of abelian
  sheaves, and the six functors are obtained from those of layer (A) with
  $\mathcal{A} = \Ab$ by transporting module structures and changing scalars. This is the
  folder \lean{SixFunctors/ProEtale/ModuleObject}.
\end{enumerate}

Layer (B) is the one that matches Bhatt--Scholze for $\ell$-adic coefficients: an
$\ell$-adic sheaf there is a sheaf of modules over the sheaf of rings $U \mapsto C(U,
\mathbb{Z}_\ell)$, which is \emph{not} the constant sheaf with value $\mathbb{Z}_\ell$
(the latter would be the sheaf of locally constant functions, which is wrong on
w-contractible objects).

\paragraph{Conventions.} Schemes live in \texttt{Type u}. The small pro-\'etale site
$X_\proet$ is mathlib's \lean{Scheme.ProEt}: schemes \emph{weakly \'etale} over $X$, with
the topology \lean{Scheme.ProEt.topology} generated by quasi-compact (fpqc) jointly
surjective families of weakly \'etale maps. Its objects are in \texttt{Type (u+1)} and
its morphisms in \texttt{Type u}; it is essentially small in \texttt{u+1}. Coefficient
categories have $\Hom$-types in \texttt{Type (u+1)}. Note that, unlike the \'etale site,
joint surjectivity alone does \emph{not} give a pro-\'etale cover: quasi-compactness is part
of the covering condition, which is why Zariski covers are fed in explicitly
(\lean{Scheme.zariskiPrecoverage_le_proetalePrecoverage}).

%======================================================================
\section{Monoidal structure on pro-\'etale sheaves}
%======================================================================

\begin{choice}[Tensor product of sheaves]
For a closed braided (resp.\ symmetric) monoidal coefficient category $\mathcal{A}$, the
tensor product on $\Sh(X_\proet,\mathcal{A})$ is the \emph{sheafification of the
objectwise tensor product of presheaves} (mathlib's \lean{Sheaf.monoidalCategory}),
installed as a global instance (\lean{ProEtale.sheafMonoidalCategory}).
\end{choice}

Closedness of $\mathcal{A}$ is what makes the class of sheafification-inverting morphisms
stable under tensoring (tensoring is a left adjoint whose right adjoint, the internal hom,
preserves sheaves). The sheaf category is then itself monoidal closed by \emph{Day's
reflection theorem} applied to $\mathrm{sheafify} \dashv \mathrm{forget}$, so that
$\mathcal{F}\otimes -$ preserves all colimits. This is the input needed for the K-flat
derivation of the tensor product (\S\ref{sec:derived}).

For $\mathcal{A} = \Ab$, the relevant monoidal structure on $\Ab$ is the tensor product
$M\otimes_{\mathbb{Z}}N$ with unit $\mathbb{Z}$ (not the cartesian one); it is
constructed on \texttt{ring-objects-for-coefficients} (\lean{AddCommGrpCat.monoidalCategory}),
together with its symmetry and closedness (internal hom $\Hom(M,N)$).

%======================================================================
\section{Coefficient rings as ring objects}
\label{sec:coefficients}
%======================================================================

\subsection{Ring objects}

A \emph{ring object} of a monoidal category $\mathcal{C}$ is a monoid object
$(R,\mu,\eta)$ (mathlib's \lean{MonObj}), commutative if $\mu\circ\beta_{R,R}=\mu$
(\lean{IsCommMonObj}). When $\mathcal{C}$ is monoidal preadditive, the global sections
$\Gamma(R)=\Hom(\mathbb{1},R)$ form a ring for the convolution product.

\begin{choice}[Rings as monoid objects of $\Ab$]
A ring $R$ is regarded as a monoid object of $(\Ab,\otimes_{\mathbb{Z}})$. Consequently:
\begin{itemize}
\item a presheaf of rings $P\colon\mathcal{C}^{op}\to\mathrm{Ring}$ is a ring object of
  abelian presheaves, objectwise;
\item since sheafification is monoidal, the sheafification of a presheaf of rings is a
  ring object of abelian sheaves (\lean{Sheaf.sheafifyRingPresheaf});
\item the abelian sheaf underlying a sheaf of rings $S$ is a ring object, the structure
  being transported along $\mathrm{sheafify}(S) \cong S$
  (\lean{Sheaf.monObj_sheafCompose_obj}); morphisms of sheaves of rings are morphisms of
  ring objects.
\end{itemize}
\end{choice}

\subsection{The $\ell$-adic coefficient sheaves}

\begin{choice}[Sheaves of continuous functions]
For a topological ring $T$, the coefficient sheaf is
\[
  \mathcal{A}_X(T) \;:=\; \bigl(U \mapsto C(U,T)\bigr)
  \quad\text{on } X_\proet,
\]
a sheaf of rings (it is even an fpqc sheaf; \lean{ProEtale.continuousMapRingSheaf}),
whose underlying abelian sheaf is a commutative ring object of $\Sh(X_\proet,\Ab)$
(\lean{ProEtale.continuousMapRingObj}). Then
\begin{align*}
\mathbb{Z}_{\ell,X} &= \mathcal{A}_X(\mathbb{Z}_\ell),&
\mathbb{Q}_{\ell,X} &= \mathcal{A}_X(\mathbb{Q}_\ell),\\
E_X &= \mathcal{A}_X(E),&
\cO_{E,X} &= \mathcal{A}_X(\cO_E),
\end{align*}
for a subfield $E\subseteq\Qlbar$ (given by an embedding $E\hookrightarrow\Qlbar$) with
the topology restricted from $\Qlbar$, and $\cO_E$ its ring of integers
(\lean{ProEtale.padicIntSheaf}, \lean{ProEtale.padicSheaf},
\lean{ProEtale.padicSubfieldSheaf}, \lean{ProEtale.padicIntegersSheaf}). The abelian sheaf
underlying $\mathbb{Z}_{\ell,X}$ is definitionally mathlib's
\lean{Scheme.ellAdicSheaf}.
\end{choice}

No completeness, Noetherianity or inverse-limit presentation of $T$ is used: the same
construction covers $\mathbb{Z}_\ell$, $\mathbb{Q}_\ell$, $E$ and $\cO_E$ uniformly. In
particular $\mathbb{Q}_{\ell,X}$ is \emph{not} defined as $\mathbb{Z}_{\ell,X}[1/\ell]$
and $E_X$ is not defined as $\cO_{E,X}[1/\ell]$; these identifications are theorems to be
proved, not definitions.

\begin{choice}[$\Qlbar$]
$\Qlbar{}_{,X}$ is \emph{not} $C(-,\Qlbar)$ for the valuation topology. Over a
w-contractible $U$ the correct sections are the continuous maps $U\to\Qlbar$ with values
in some \emph{finite} subextension $E\subseteq\Qlbar$ of $\mathbb{Q}_\ell$; this is
strictly smaller than $C(U,\Qlbar)$ (a null sequence of elements of unbounded degree
extends to $\beta\mathbb{N}$). Accordingly
\[
  \Qlbar{}_{,X} := \mathrm{sheafify}\Bigl(U\mapsto \{f\in C(U,\Qlbar) \mid
  f(U)\subseteq E \text{ for some finite } E/\mathbb{Q}_\ell\}\Bigr)
\]
(\lean{ProEtale.padicAlgClSheaf}), defined directly as a ring object of abelian sheaves.
This is $\operatorname{colim}_E E_X$ of Bhatt--Scholze \S6.8.
\end{choice}

The inclusions $\mathbb{Z}_{\ell,X}\to\mathbb{Q}_{\ell,X}\to E_X\to\Qlbar{}_{,X}$,
$\mathbb{Z}_{\ell,X}\to\cO_{E,X}\to E_X$ are morphisms of ring objects.

\subsection{Coefficient morphisms along a morphism of schemes}

For $f\colon X\to Y$ the rings $\mathcal{A}_X(T)$ and $\mathcal{A}_Y(T)$ are related by
two canonical morphisms of ring objects (\lean{ProEtale.continuousMapPushforwardHom},
\lean{ProEtale.continuousMapPullbackHom}):
\[
  \eta_f\colon \mathcal{A}_Y(T)\to f_*\mathcal{A}_X(T),\qquad
  (h\in C(V,T))\mapsto h\circ p_V,\quad p_V\colon V\times_Y X\to V,
\]
and its mate under $f^*\dashv f_*$,
\[
  \alpha_f\colon f^*\mathcal{A}_Y(T)\to\mathcal{A}_X(T).
\]
$\eta$ is compatible with identities and compositions. $\alpha_j$ is an isomorphism for an
open immersion $j$ (indeed for weakly \'etale $j$), but \emph{not} in general, e.g.\ not
for closed immersions. This is the reason the module-object pullback below needs an
extension of scalars.

%======================================================================
\section{Module objects and change of scalars}
\label{sec:modules}
%======================================================================

\begin{choice}[Sheaves of modules]
A sheaf of $\mathcal{R}$-modules, for a ring object $\mathcal{R}$ of
$\Sh(X_\proet,\mathcal{A})$, is a left \emph{module object} $(M,\gamma\colon
\mathcal{R}\otimes M\to M)$ (mathlib's \lean{ModObj}, category \lean{Mod C R}). The
category of $\ell$-adic sheaves is
\[
  \Mod_{\mathcal{A}_X(T)} := \Mod\bigl(\Sh(X_\proet,\Ab),\,\mathcal{A}_X(T)\bigr).
\]
\end{choice}

This choice avoids a separate theory of sheaves of modules over a varying sheaf of rings:
everything is internal to the fixed closed symmetric monoidal category
$\Sh(X_\proet,\Ab)$.

\subsection{Structure of $\Mod_{\mathcal{R}}(\mathcal{C})$}
Let $\mathcal{C}$ be abelian, monoidal preadditive, and $\mathcal{R}$ a ring object with
$\mathcal{R}\otimes-$ right exact.
\begin{itemize}
\item $\Mod_{\mathcal{R}}(\mathcal{C})$ is abelian, with kernels, cokernels and biproducts
  computed in $\mathcal{C}$; the forgetful functor is exact and faithful.
\item For $\mathcal{C}$ braided with coequalizers preserved by $\otimes$ and $\mathcal{R}$
  commutative, the \emph{relative tensor product}
  $M\otimes_{\mathcal{R}}N=\operatorname{coeq}(M\otimes\mathcal{R}\otimes N\rightrightarrows
  M\otimes N)$ makes $\Mod_{\mathcal{R}}$ monoidal with unit $\mathcal{R}$; it is obtained
  by identifying $\mathcal{R}$-modules with symmetric $\mathcal{R}$-$\mathcal{R}$-bimodules
  and inducing from mathlib's tensor product of bimodules. If $\mathcal{C}$ is symmetric,
  so is $\Mod_{\mathcal{R}}$.
\item For $\mathcal{C}$ closed symmetric Grothendieck abelian, $\Mod_{\mathcal{R}}$ is
  again closed (internal hom $\Hom_{\mathcal{R}}(M,N)$ as an equalizer of
  $[M,N]\rightrightarrows[\mathcal{R}\otimes M,N]$), Grothendieck abelian (separator
  $\mathcal{R}\otimes G$), and inherits functorial flat resolutions
  ($\mathcal{R}\otimes P(M)\to M$).
\end{itemize}

\subsection{Change of scalars}
For a morphism of ring objects $\varphi\colon\mathcal{R}\to\mathcal{S}$
(\lean{CategoryTheory/ModuleObjects/ExtendScalars}):
\begin{itemize}
\item \emph{restriction} $\Res_\varphi\colon\Mod_{\mathcal{S}}\to\Mod_{\mathcal{R}}$
  (mathlib's \lean{Mod.comap}), exact;
\item \emph{extension} $\Ext_\varphi = \mathcal{S}\otimes_{\mathcal{R}}-$, for
  $\mathcal{R}$ commutative, built as the bimodule tensor product of $\mathcal{S}$ (an
  $\mathcal{S}$-$\mathcal{R}$-bimodule via $\varphi$) with $M$; right exact, not exact;
\item \emph{coextension} $\Coext_\varphi = \Hom_{\mathcal{R}}(\mathcal{S},-)$, for
  $\mathcal{C}$ closed, an equalizer of $[\mathcal{S},N]\rightrightarrows
  [\mathcal{R}\otimes\mathcal{S},N]$;
\end{itemize}
with adjunctions $\Ext_\varphi\dashv\Res_\varphi\dashv\Coext_\varphi$, and compatibilities
with identities and compositions. For $\varphi$ an isomorphism,
$\Ext_\varphi\cong\Res_{\varphi^{-1}}\cong\Coext_\varphi$.

\paragraph{Derived extension of scalars.} $\mathcal{S}\lotimes_{\mathcal{R}}- :=
L(\Ext_\varphi)$, the total left derived functor computed on K-flat resolutions. Since
$\Res_\varphi$ is exact, $R(\Res_\varphi)=\D(\Res_\varphi)$ is levelwise, and
$\mathcal{S}\lotimes_{\mathcal{R}}-\dashv\D(\Res_\varphi)$.

\subsection{Transport along monoidal functors and projection formulas}
\begin{itemize}
\item A lax monoidal functor $F\colon\mathcal{C}\to\mathcal{D}$ lifts to
  $\mathrm{mapMod}(F)\colon\Mod_{\mathcal{R}}(\mathcal{C})\to
  \Mod_{F\mathcal{R}}(\mathcal{D})$, with action $F\mathcal{R}\otimes FM\to
  F(\mathcal{R}\otimes M)\to FM$. It is exact when $F$ is.
\item A \emph{monoidal adjunction} $F\dashv G$ ($F$ strong, $G$ lax, structures mates)
  lifts to $\mathrm{mapMod}(F)\dashv\Res_{\mathrm{unit}}\circ\mathrm{mapMod}(G)$.
\item \emph{Projection formula.} For $S$ strong monoidal and $P$ a functor, a projection
  formula is a coherent natural isomorphism $P(SL\otimes M)\cong L\otimes PM$. When $P\dashv
  S$ the canonical map is the adjoint of $SL\otimes M\to SL\otimes SPM\to S(L\otimes PM)$,
  and the projection formula is the condition that it be invertible. A projection formula
  \emph{linearizes} $P$ to $P_B\colon\Mod_{SB}\to\Mod_B$, with action
  $B\otimes PM\cong P(SB\otimes M)\to PM$, and any right adjoint $G$ of $P$ to $G_B$, with
  $P_B\dashv G_B$.
\end{itemize}
These are the formal tools from which all module-object six functors are assembled.

%======================================================================
\section{Derived categories}
\label{sec:derived}
%======================================================================

\begin{choice}[Unbounded derived categories, K-flat / K-injective derivation]
All derived categories are mathlib's \lean{DerivedCategory}, the localization of cochain
complexes at quasi-isomorphisms, for the Grothendieck abelian categories at hand. Right
derived functors are total derived functors computed on K-injective resolutions; left
derived functors (pullback, tensor product, extension of scalars) are total derived
functors computed on K-flat resolutions, via a left derivability structure that requires a
monoidal preadditive structure whose tensor preserves colimits and functorial flat
resolutions. Bounded below versions use \lean{DerivedCategory.Plus}.
\end{choice}

The derived categories present in the development:

\medskip
\noindent
\begin{tabularx}{\textwidth}{@{}lX@{}}
\toprule
Category & Definition \\
\midrule
$\D(X_\proet,\mathcal{A})$ &
  $\lean{DerivedCategory}(\Sh(X_\proet,\mathcal{A}))$, fixed coefficients (layer A);
  $\D^+$ its bounded below part. \\
$\D(\Mod_{\mathcal{A}_X(T)})$ &
  $\lean{DerivedCategory}(\Mod(\Sh(X_\proet,\Ab),\mathcal{A}_X(T)))$, module objects
  (layer B). For $T=\cO_E$ resp.\ $E$ this is $\D(X_\proet,\cO_E)$ resp.\
  $\D(X_\proet,E)$. \\
$\D_{\mathrm{cons}}(X_\proet,\mathcal{R})$ &
  full subcategory of $\D(\Mod_{\mathcal{R}})$ of bounded complexes with constructible
  cohomology sheaves (\S\ref{sec:constructible}). \\
$\D_{\mathrm{comp}}(J,\cO_E)$ &
  (\texttt{master}) full subcategory of $\D(\Sh(J,\Mod_{\cO_E}))$ of derived
  $\mathfrak{m}_E$-complete objects: $K\to R\lim_n K\lotimes\cO_E/\mathfrak{m}_E^n$ is an
  isomorphism. \\
$\D_{\mathrm{comp}}(J,E)$ &
  (\texttt{master}) Verdier quotient of $\D_{\mathrm{comp}}(J,\cO_E)$ by objects killed by
  a nonzero element of $\cO_E$; $E$-linear triangulated. \\
\bottomrule
\end{tabularx}
\medskip

\begin{remark}[Two models of $\ell$-adic coefficients]
The derived-complete categories on \texttt{master} use \emph{constant} coefficients:
sheaves with values in $\Mod_{\cO_E}$ on an abstract site $J$, made complete afterwards.
The new layer (B) uses sheaves of modules over the \emph{non-constant} ring sheaf
$\cO_{E,X}=C(-,\cO_E)$, as in Bhatt--Scholze \S6.8; there $E$-coefficients are modules over
$E_X=C(-,E)$ rather than a Verdier quotient. No derived completeness is imposed in layer
(B): completeness is regarded as a choice of subcategory on which the functors act, and
does not affect the coefficient morphisms or the constructions of the functors. Relating
the two models is not done.
\end{remark}

\paragraph{Monoidal structure.} On $\D(X_\proet,\mathcal{A})$ and on
$\D(\Mod_{\mathcal{R}})$ the tensor product is the derived tensor product $\lotimes$
(resp.\ $\lotimes_{\mathcal{R}}$), the two-variable total left derived functor of $\otimes$
(resp.\ the relative tensor product $\otimes_{\mathcal{R}}$), making the localization
functor lax monoidal. The derived internal hom $\RHom$ is \emph{defined} as the internal
hom of the resulting closed monoidal structure, i.e.\ as the right adjoint of
$\lotimes$: $K\lotimes-\dashv\RHom(K,-)$.

%======================================================================
\section{Constructibility}
\label{sec:constructible}
%======================================================================

From the branch \texttt{bounded-constructible-module-object-of-sheaf}
(\lean{ConstructibleSheaf/ProEtaleModule}). Let $\mathcal{R}$ be a ring object of
$\Sh(X_\proet,\mathcal{A})$, $\mathcal{A}$ closed symmetric Grothendieck abelian. For a
morphism $g\colon Z\to X$, ``restriction'' of an $\mathcal{R}$-module $M$ along $g$ means
the $g^*\mathcal{R}$-module $\mathrm{mapMod}(g^*)(M)=g^*M$, using that $g^*$ is monoidal.

\begin{definition}[Free and finitely presented]
$\mathcal{R}^n:=\bigoplus_{i<n}\mathcal{R}$. A module is \emph{finitely presented} if it
is isomorphic to $\operatorname{coker}(\mathcal{R}^m\to\mathcal{R}^n)$.
\end{definition}

\begin{definition}[Locally finitely presented]
$M\in\Mod_{\mathcal{R}}$ is \emph{locally finitely presented}
(\lean{ProEtale.IsLocallyFinitelyPresented}) if there is a pro-\'etale cover
$\{U_i\to X\}$ such that each $M|_{U_i}$ is a finitely presented
$\mathcal{R}|_{U_i}$-module. Torsion is allowed.
\end{definition}

\begin{definition}[Constructible, Bhatt--Scholze \S6.8]
$M$ is \emph{constructible} (\lean{ProEtale.IsConstructible}) if there is a finite
stratification $X=\coprod_a X_a$ by locally closed subschemes
(\lean{Scheme.Stratification}, the same notion as for \'etale constructible sheaves) such
that each $M|_{X_a}$ is locally finitely presented over $\mathcal{R}|_{X_a}$.
\end{definition}

For $\mathcal{R}=\cO_{E,X}$ these are the constructible $\cO_E$-sheaves; for
$\mathcal{R}=E_X$ the constructible $E$-sheaves (lisse on strata).

\begin{definition}[Bounded constructible derived category]
For $\mathcal{R}$ commutative,
\[
  \D_{\mathrm{cons}}(X_\proet,\mathcal{R}) \subseteq \D(\Mod_{\mathcal{R}})
\]
is the full subcategory (\lean{ProEtale.BoundedConstructible}) of complexes $K$ that are
bounded for the canonical t-structure ($H^i(K)=0$ outside a finite interval) and all of
whose cohomology sheaves $H^i(K)$ are constructible. Morphisms are computed in
$\D(\Mod_{\mathcal{R}})$; no derived completeness condition is imposed. On a
topologically Noetherian $X$ it is a triangulated subcategory (constructible modules being
closed under kernels, cokernels and extensions).
\end{definition}

\begin{remark}
The definition is made for an arbitrary ring object $\mathcal{R}$ over a fixed $X$, so it
applies directly to $\mathcal{R}=\mathcal{A}_X(T)$. Restriction to a stratum lands in
$g^*\mathcal{R}$-modules, \emph{not} in $\mathcal{A}_{X_a}(T)$-modules; these differ for a
closed stratum since $\alpha_g$ is not an isomorphism there. Whether constructibility of a
$\mathcal{A}_X(T)$-module should be tested after extension of scalars along $\alpha_g$
(i.e.\ via the module-object pullback $g^*$ of \S\ref{sec:sixB}) is a point to align when
the two branches are merged.
\end{remark}

%======================================================================
\section{Six functors with fixed coefficients (layer A)}
\label{sec:sixA}
%======================================================================

Throughout, $\mathcal{A}$ is a closed symmetric monoidal Grothendieck abelian category.

\subsection{Pushforward and pullback}
\begin{choice}
For an \emph{arbitrary} $g\colon Z\to Y$, base change
$Y_\proet\to Z_\proet$, $V\mapsto V\times_YZ$, is continuous. Then
\[
  (g_*G)(V) := G(V\times_Y Z),\qquad
  g^* := \mathrm{sheafify}\circ\operatorname{Lan}_{(\text{base change})^{op}},
\]
with $g^*\dashv g_*$ (\lean{ProEtale.lowerStar}, \lean{ProEtale.upperStar}).
\end{choice}
Defining both along base change (rather than along the localization functor
$Z_\proet\to Y_\proet$, $W\mapsto(W\to Z\to Y)$, which needs $g$ weakly \'etale) is essential
for compactifications, whose proper part is not weakly \'etale. For $g$ weakly \'etale,
$g^*$ is restriction along the localization functor. Properties: $g_*$ is left exact and
lax symmetric monoidal; $g^*$ is exact, and strong symmetric monoidal with
$g^*\dashv g_*$ a monoidal adjunction (the oplax structure of $g^*$ is the mate of the lax
structure of $g_*$); both compose functorially.

\subsection{Extension by zero and lower shriek}
\begin{choice}
For an open immersion $j\colon U\to X$, $j_!:=\mathrm{sheafify}\circ
\operatorname{Lan}_{(\text{localization})^{op}}$, left adjoint to restriction
$j^*$ (\lean{openImmersionLowerShriek}). For $f\colon X\to Y$ with a chosen
compactification $X\xrightarrow{i}\bar X\xrightarrow{g}Y$ ($i$ open immersion, $g$
proper),
\[
  f_! := g_*\circ i_!
\]
(\lean{ProEtale.lowerShriek}), definitionally attached to the chosen compactification
\lean{HasCompactification.compactification}.
\end{choice}
Properties: $j_!\dashv j^*\dashv j_*$ is an adjoint triple; $j_!$ is exact and fully
faithful; $f_!$ is left exact. There is a canonical comparison $f_!\to f_*$ (from
$i_!\to i_*$), which is an isomorphism for $f$ proper (proved via the vanishing of both
sides off the range of $i$ and conservativity of restriction to the Zariski cover
$\{i,\ \text{complement}\}$).

\subsection{Derived functors}
\begin{itemize}
\item $Rg_* := R(g_*)$ (K-injective), also on $\D^+$.
\item $Lg^* := L(g^*)$ (K-flat). Since $g^*$ is exact, $Lg^*\cong\D(g^*)$ is the levelwise
  functor, compatible with composition, with $H^n$ and with degree-$0$ objects; one has
  $Lg^*\dashv Rg_*$ on $\D$ and $\D^+(g^*)\dashv R^+g_*$ on $\D^+$.
\item For an open immersion, $\D(j_!)\dashv\D(j^*)$ (both exact), also on $\D^+$.
\item \textbf{Derived lower shriek.}
  \begin{choice}
  $Rf_! := R(g_*)\circ\D(i_!)$, for the chosen compactification; on $\D^+$,
  $R^+(g_*)\circ\D^+(i_!)$, which is the restriction of the unbounded one.
  \end{choice}
  It is deliberately \emph{not} the total derived functor of $f_!=g_*i_!$: $i_!$ does not
  send injectives to $g_*$-acyclics, and $R(g_*i_!)$ is the wrong functor, although
  $H^0(Rf_!)=f_!$.
\item \textbf{Exceptional pullback.}
  \begin{choice}
  $Rf^! := \D(i^*)\circ Rg^!$, where $Rg^!$ is a right adjoint of $Rg_*$ for the proper
  $g$ (which exists by Brown representability; it is an abstract choice of right adjoint,
  used only through its adjunction).
  \end{choice}
  Then $Rf_!\dashv Rf^!$ is the composite of $\D(i_!)\dashv\D(i^*)$ and $Rg_*\dashv Rg^!$;
  in particular, that $Rf_!$ is a left adjoint is a theorem, not an assumption. On $\D^+$,
  $Rf^!$ is defined as the restriction of the unbounded $Rf^!$ (which preserves bounded
  below objects), and the adjunction restricts.
\item \textbf{Tensor and hom.} $\lotimes$ and $\RHom$ on $\D(X_\proet,\mathcal{A})$ as in
  \S\ref{sec:derived}.
\item \textbf{Cohomology.} $R\Gamma(X,-):=R\Gamma$ of global sections;
  $R\Gamma_c(X/Y,-):=R\Gamma(Y,Rf_!-)$; $H^i_c:=H^i\circ R\Gamma_c$.
\end{itemize}

%======================================================================
\section{Six functors on module objects (layer B)}
\label{sec:sixB}
%======================================================================

Let $T$ be a commutative topological ring and write $\mathcal{A}_X=\mathcal{A}_X(T)$. All
functors are of the form ``transport along the $\Ab$-valued functor, then change scalars
along $\eta_f$ or $\alpha_f$''. Nothing requires $\alpha_f$ to be invertible except the
open-immersion statements.

\subsection{Pushforward and pullback}
\begin{choice}
\begin{align*}
  f_* &:= \Res_{\eta_f}\circ\mathrm{mapMod}(f_*)\colon
     \Mod_{\mathcal{A}_X}\to\Mod_{f_*\mathcal{A}_X}\to\Mod_{\mathcal{A}_Y},\\
  f^* &:= \Ext_{\alpha_f}\circ\mathrm{mapMod}(f^*)\colon
     \Mod_{\mathcal{A}_Y}\to\Mod_{f^*\mathcal{A}_Y}\to\Mod_{\mathcal{A}_X}.
\end{align*}
\end{choice}
So the underlying sheaf of $f_*M$ is $f_*M$ with $\mathcal{A}_Y$ acting through $\eta_f$,
and $f^*M=\mathcal{A}_X\otimes_{f^*\mathcal{A}_Y}f^*M$ is the ringed-topos pullback.
$f_*$ is left exact; $f^*$ is right exact but \emph{not} exact in general.
The adjunction $f^*\dashv f_*$ is the composite of the lifted monoidal adjunction
$\mathrm{mapMod}(f^*)\dashv\Res_{\mathrm{unit}}\circ\mathrm{mapMod}(f_*)$ and
$\Ext_{\alpha_f}\dashv\Res_{\alpha_f}$, using that $\eta_f$ factors as the unit followed by
$f_*\alpha_f$. One also has the \emph{linear pushforward}
$\Res_{\mathrm{unit}}\circ\mathrm{mapMod}(f_*)\colon\Mod_{f^*\mathcal{A}_Y}\to
\Mod_{\mathcal{A}_Y}$, with $f_*=(\text{linear }f_*)\circ\Res_{\alpha_f}$.

\subsection{Extension by zero and lower shriek}
\begin{choice}
For an open immersion $j\colon U\to X$, the $\Ab$-valued $j_!$ satisfies the projection
formula $j_!(j^*L\otimes M)\cong L\otimes j_!M$; this linearizes it to
$\text{linear }j_!\colon\Mod_{j^*\mathcal{A}_X}\to\Mod_{\mathcal{A}_X}$, and
\[
  j_! := (\text{linear }j_!)\circ\Res_{\alpha_j}\colon
  \Mod_{\mathcal{A}_U}\to\Mod_{\mathcal{A}_X}.
\]
For $f$ with chosen compactification $(i,g)$: $f_!:=g_*\circ i_!$ on module objects.
\end{choice}
The underlying sheaf of $j_!M$ (resp.\ $f_!M$) is the $\Ab$-valued $j_!M$ (resp.\ $f_!M$).
$j_!$ is exact and, since $\alpha_j$ is an isomorphism, $j^*$ is exact and
$j_!\dashv j^*$ (composite of $\text{linear }j_!\dashv\mathrm{mapMod}(j^*)$ with
$\Res_{\alpha_j}\dashv\Coext_{\alpha_j}\cong\Ext_{\alpha_j}$). $f_!$ is left exact. The
projection-formula route is the template intended for $f^!$ as well:
$f_!=P_B\circ\Res_{\alpha_f}$, $f^!=\Coext_{\alpha_f}\circ G_B$.

\subsection{Derived functors}
\begin{itemize}
\item $Rf_* := R(f_*)$ (K-injective in $\Mod_{\mathcal{A}_X}$), also on $\D^+$. Since
  $\Res_{\eta_f}$ is exact,
  $Rf_*\cong\D(\Res_{\eta_f})\circ R(\mathrm{mapMod}(f_*))$.
\item $Lf^* := L(f^*)$ (K-flat in $\Mod_{\mathcal{A}_Y}$). Since $\mathrm{mapMod}(f^*)$ is
  exact and preserves K-flat complexes,
  \[
    Lf^*M \;\cong\; \mathcal{A}_X\lotimes_{f^*\mathcal{A}_Y} f^*M,
    \qquad Lf^*\cong L(\Ext_{\alpha_f})\circ\D(\mathrm{mapMod}(f^*)),
  \]
  the pullback of Bhatt--Scholze, Remark 6.8.15. Unlike layer (A), $Lf^*$ is not the
  levelwise functor. $Lf^*\dashv Rf_*$ on $\D$; there is no general $\D^+$ adjunction.
\item For an open immersion, $Lj^*\cong\D(j^*)$ and $\D(j_!)\dashv\D(j^*)$, also on
  $\D^+$.
\item $Rf_! := R(g_*)\circ\D(i_!)$ for the chosen compactification, and on $\D^+$, as in
  layer (A).
\item $\lotimes_{\mathcal{A}_X}$ and $\RHom_{\mathcal{A}_X}$ on $\D(\Mod_{\mathcal{A}_X})$
  as in \S\ref{sec:derived}.
\item $Rf^!$ on module objects is \emph{not yet constructed}.
\end{itemize}

%======================================================================
\section{Summary}
%======================================================================

\medskip
\noindent
\begin{tabularx}{\textwidth}{@{}lXX@{}}
\toprule
Functor & Layer A: $\Sh(X_\proet,\mathcal{A})$ & Layer B: $\Mod_{\mathcal{A}_X(T)}$\\
\midrule
$f_*$ & restriction along base change &
  $\Res_{\eta_f}\circ\mathrm{mapMod}(f_*)$\\
$f^*$ & $\mathrm{sheafify}\circ\operatorname{Lan}$, exact, monoidal &
  $\Ext_{\alpha_f}\circ\mathrm{mapMod}(f^*)$, right exact\\
$j_!$ ($j$ open) & $\mathrm{sheafify}\circ\operatorname{Lan}$ along localization &
  $(\text{linear }j_!)\circ\Res_{\alpha_j}$ via projection formula\\
$f_!$ & $g_*\circ i_!$ (chosen compactification) & $g_*\circ i_!$\\
$Rf_*$ & $R(f_*)$ & $R(f_*)\cong\D(\Res_{\eta_f})\circ R(\mathrm{mapMod}f_*)$\\
$Lf^*$ & $L(f^*)\cong\D(f^*)$ & $L(f^*)\cong\mathcal{A}_X\lotimes_{f^*\mathcal{A}_Y}f^*$\\
$Rf_!$ & $R(g_*)\circ\D(i_!)$ & $R(g_*)\circ\D(i_!)$\\
$Rf^!$ & $\D(i^*)\circ Rg^!$, $Rg^!$ right adjoint of $Rg_*$ & not yet\\
$\lotimes$ & K-flat derived $\otimes$ & K-flat derived $\otimes_{\mathcal{A}_X}$\\
$\RHom$ & right adjoint of $\lotimes$ & right adjoint of $\lotimes_{\mathcal{A}_X}$\\
\bottomrule
\end{tabularx}
\medskip

\paragraph{Main open directions.}
\begin{itemize}
\item $Rf^!$ on module objects, via the projection-formula linearization of a right
  adjoint and coextension of scalars along $\alpha_f$.
\item Coefficients $\Qlbar{}_{,X}$, which are not of the form $\mathcal{A}_X(T)$.
\item Restricting the module-object functors to $\D_{\mathrm{cons}}$ (and to derived
  complete objects for $\cO_E$), and aligning the restriction-to-strata used in
  constructibility with the module-object pullback.
\item Comparing layer (B) with the derived-complete constant-coefficient model on
  \texttt{master}, and identifying $E_X$ with $\cO_{E,X}[1/\ell]$.
\end{itemize}

\paragraph{Branches.}
\texttt{ring-objects-for-coefficients} (\S\ref{sec:coefficients}) and
\texttt{extend-scalars-module-object} (\S\ref{sec:modules}) are merged into
\texttt{module-object-six-functors} (\S\ref{sec:sixB}, and the move of layer (A) to
\lean{FiniteCoefficient}). \texttt{bounded-constructible-module-object-of-sheaf}
(\S\ref{sec:constructible}) is not yet merged with it.

\end{document}
